\documentclass[11pt,class=report,crop=false]{standalone}
\usepackage[screen]{../logic}


\begin{document}


%====================================================================
\chapitre{Théorème de complétude}
%====================================================================

% \insertvideo{}{partie 1.1. Titre}


\objectifs{Le théorème de complétude dit que, en logique du premier ordre, tout ce qui est vrai est prouvable. Nous énonçons ce théorème, ses variantes et ses applications.}


%%%%%%%%%%%%%%%%%%%%%%%%%%%%%%%%%%%%%%%%%%%%%%%%%%%%%%%%%%%%%%%%%%%%%
\section{Théorème de complétude}

%--------------------------------------------------------------------
\subsection{Énoncé}

Soit $\mathcal{L}$ un langage du premier ordre.

\begin{theoreme}[Théorème de complétude]
	\sauteligne
\mybox{
	Si \quad 
	$\mathcal{P}_1, \ldots, \mathcal{P}_l \quad \ldoubleturn \quad \mathcal{Q}$
	\quad \text{ alors }
	\quad 
	$\mathcal{P}_1, \ldots, \mathcal{P}_l \quad \lturn \quad \mathcal{Q}$
}	
\end{theoreme}
\index{theoreme completude@théorème de complétude}


Ce que l'on peut lire : si, lorsque les hypothèses $\mathcal{P}_1,\ldots,\mathcal{P}_l$ sont satisfaites, $\mathcal{Q}$ l'est aussi, alors il existe une preuve de $\mathcal{Q}$ à partir de $\mathcal{P}_1,\ldots,\mathcal{P}_l$. 

Voici la variante pour un ensemble $\Gamma$ de formules, éventuellement infini dénombrable :
\begin{theoreme}[Théorème de complétude]
	\sauteligne
	\[
		\text{Si} \quad 
		\Gamma \quad \ldoubleturn \quad \mathcal{Q}
		\quad \text{ alors }
		\quad 
		\Gamma \quad \lturn \quad \mathcal{Q}.
	\]	
\end{theoreme}

La conclusion $\Gamma \ \lturn \ \mathcal{Q}$ est un résultat d'existence : \og{}il existe une preuve de $\mathcal{Q}$ à partir des formules de $\Gamma$\fg{}. Cependant, le théorème n'est pas constructif et ne donne aucune indication concrète sur cette preuve.

La démonstration du théorème de complétude est assez exigeante. Elle partage des points communs avec le théorème de compacité de la logique V/F mais avec un niveau de difficulté supplémentaire.
Cette démonstration sera faite au prochain chapitre. On va maintenant se concentrer sur les différentes reformulations et les applications.

%--------------------------------------------------------------------
\subsection{Équivalence}

On a vu au chapitre précédent le théorème de validité qui est exactement la réciproque du théorème de complétude.
Ainsi, on obtient une équivalence qu'on appelle encore théorème de complétude :

\begin{theoreme}[Théorème de complétude]
	\sauteligne
\mybox{$	
	\Gamma \quad \ldoubleturn \quad \mathcal{Q}
	\quad \text{ ssi }
	\quad 
	\Gamma \quad \lturn \quad \mathcal{Q}
$}
\end{theoreme}

Ainsi, pour la logique du premier ordre, la conséquence logique et la prouvabilité sont des notions équivalentes.



%--------------------------------------------------------------------
\subsection{Et l'incomplétude ?}

N'y a-t-il pas une contradiction à obtenir le théorème de complétude alors qu'à la fin du livre on obtiendra le théorème d'incomplétude, dont le nom semble suggérer le contraire ?
En fait, le théorème de complétude s'applique à la logique du premier ordre qui est une logique assez rudimentaire mais avec l'hypothèse $\Gamma \ldoubleturn \mathcal{Q}$ qui est assez forte car 
la conséquence logique exige que la formule soit vraie dans toutes les structures possibles.

Le théorème d'incomplétude ne concerne pas l'ensemble des structures, mais concerne des théories plus spécifiques, typiquement qui contiennent les entiers naturels $\Nn$ avec les axiomes de l'arithmétique (les axiomes de Peano).
On verra que, pour l'arithmétique des entiers, il existe des formules $\mathcal{Q}$ telles que ni $\mathcal{Q}$ ni $\lnot \mathcal{Q}$ ne sont prouvables, autrement dit, il existe des formules vraies mais non prouvables, mais tout ceci uniquement pour la structure associée à l'arithmétique des entiers.
Ceci n'est pas en contradiction avec le théorème de complétude car l'hypothèse \og{}$\Gamma \ldoubleturn_{\mathcal{S}} \mathcal{Q}$\fg{} n'est pas vérifiée pour toutes les structures, mais seulement pour certaines structures spécifiques ; les conditions du théorème de complétude ne sont donc pas réunies. 

Ainsi, les théorèmes de complétude et d'incomplétude sont tous les deux vrais, mais ne s'appliquent pas aux mêmes situations et ne se contredisent donc pas.


%%%%%%%%%%%%%%%%%%%%%%%%%%%%%%%%%%%%%%%%%%%%%%%%%%%%%%%%%%%%%%%%%%%%%
\section{Applications aux théories}

On va rappeler des définitions et introduire de nouvelles notions.
Soit $\mathcal{L}$ un langage du premier ordre. 

%--------------------------------------------------------------------
\subsection{Théorie cohérente, modèle}

\begin{itemize}
	\item Une \defi{formule fermée}\index{formule!fermee@fermée} (ou \defi{formule close}, ou \defi{phrase}) est une formule sans variable libre, c'est-à-dire que toutes les occurrences de variables sont dans la portée d'un quantificateur.
	
	\item Une \defi{théorie}\index{theorie@théorie} $\mathcal{T}$ est un ensemble de formules fermées. (Remarque : certains appellent \og{}théorie\fg{} un ensemble quelconque de formules.)
	
	\item Une structure $\mathcal{S}$ est un \defi{modèle}\index{modele@modèle} pour $\mathcal{T}$ si $\mathcal{T}$ est satisfaite dans $\mathcal{S}$, noté $\ldoubleturn_{\mathcal{S}} \mathcal{T}$, c'est-à-dire que l'interprétation de chaque formule de $\mathcal{T}$ est vraie dans $\mathcal{S}$. (Comme les formules sont fermées, il n'y a pas d'affectation à considérer.)
\end{itemize}

Voici une nouvelle notion très importante.

\begin{definition}
	Une théorie $\mathcal{T}$ est \defi{non cohérente} s'il existe une formule $\mathcal{Q}$ telle que :
	\[
	\mathcal{T} \ \lturn \ \mathcal{Q} \quad \text{ et } \quad \mathcal{T} \ \lturn \ \lnot \mathcal{Q}.
	\] 
	
	Dans le cas contraire, on dit que $\mathcal{T}$ est \defi{cohérente}\index{theorie@théorie!coherente@cohérente}. 
\end{definition}

Il est équivalent de dire : 
\mybox{
	$\mathcal{T}$ est non cohérente \ ssi \ $\mathcal{T} \ \lturn \ \lantitauto$
}
où $\lantitauto$ désigne une contradiction.
Ainsi, d'une théorie cohérente on ne pourra jamais démontrer une contradiction (ce qui est rassurant).



%--------------------------------------------------------------------
\subsection{Lemme de déduction naturelle}

$\Gamma$ désigne un ensemble de formules, et $\mathcal{P}$, $\mathcal{Q}$ des formules.
\begin{lemme}[Lemme de déduction naturelle] 
	\sauteligne
	\mybox{
	$\Gamma, \ \mathcal{P} \quad \lturn \quad \mathcal{Q}$ \qquad ssi \qquad $\Gamma \quad \lturn \quad (\mathcal{P} \limplies \mathcal{Q})$
}
\end{lemme}
\index{lemme!de deduction naturelle@de déduction naturelle}

\begin{proof}
Le sens réciproque $\Leftarrow$ est en fait la règle du \emph{modus ponens} $\mathcal{P}, \mathcal{P} \limplies \mathcal{Q} \ \ \lturn \ \ \mathcal{Q}$. 

Le sens direct $\Rightarrow$ est immédiat quand on a une preuve en déduction naturelle (d'où le nom du lemme), il s'agit juste d'une reformulation dans l'écriture de la preuve :	
	

\begin{minipage}[t]{0.4\textwidth}
	\qquad\textbf{Preuve de \ \ $\Gamma, \ \mathcal{P} \quad \lturn \quad \mathcal{Q}$}

	\[
	\begin{nd}
		\hypo [~] {} {\Gamma}  \by{}{}
		\hypo [~] {} {\mathcal{P}}  \by{\text{Hypothèses}}{}
		\have [~] {} {\ \vdots} \by{\text{Étapes}}{}
		\have [\lturn] {} {\mathcal{Q}} \by{\text{Conclusion}}{}
	\end{nd}
	\]
\end{minipage}
\qquad\qquad\qquad
\begin{minipage}[t]{0.4\textwidth}
	\qquad\textbf{Preuve de \ \ $\Gamma \quad \lturn \quad (\mathcal{P} \limplies \mathcal{Q})$}

	\[	
	\begin{nd}
		\hypo [~] {} {\Gamma}  \by{}{}
		\open				
		\hypo [~] {} {\mathcal{P}} \by{\text{Sous-hypothèse}}{}
		\have [~] {} {\ \vdots} \by{\text{Mêmes étapes}}{}
		\have [~] {} {\mathcal{Q}} \by{\text{Sous-conclusion}}{}
		\close 
		\have [\lturn] {} {\mathcal{P} \limplies \mathcal{Q}} \by{\text{Conclusion}}{}
	\end{nd}
	\]
\end{minipage}

\medskip	
\end{proof}


Soit $\mathcal{T}$ une théorie et $\mathcal{Q}$ une formule fermée.
\begin{corollaire}
	\sauteligne
	\mybox{$
	\mathcal{T} \quad \lturn \quad \mathcal{Q} \qquad \text{ssi} \qquad  \mathcal{T} \cup \{\lnot \mathcal{Q}\} \ \text{ non cohérente}
	$}
\end{corollaire}

\begin{proof}
\begin{align*}
					&  \mathcal{T} \cup \{\lnot \mathcal{Q}\} \ \text{ non cohérente}\\
\text{ ssi } \quad 	& \mathcal{T} \cup \{\lnot \mathcal{Q}\} \ \ \lturn \ \ \lantitauto \qquad \text{par définition}\\
\text{ ssi } \quad 	& \mathcal{T}  \ \ \lturn \ \ (\lnot\mathcal{Q} \limplies \lantitauto)  \qquad \text{par le lemme de déduction naturelle}\\
\text{ ssi } \quad 	& \mathcal{T}  \ \ \lturn \ \ \lnot(\lnot \mathcal{Q}) \qquad \text{car par définition $\lnot \mathcal{P}$ signifie $\mathcal{P} \limplies \lantitauto$}\\
\text{ ssi } \quad 	& \mathcal{T}  \ \ \lturn \ \ \mathcal{Q} \qquad \text{par double négation}
\end{align*} 	
\end{proof}

%--------------------------------------------------------------------
\subsection{Différentes formulations du théorème de complétude}

Nous allons énoncer le théorème de complétude ainsi que quelques reformulations.

Cadre : $\mathcal{T}$ est une théorie, $\mathcal{Q}$ une formule fermée.

Formulation classique :
\begin{theoreme}[Théorème A] 
	\sauteligne
\mybox{$	
	\mathcal{T} \quad \ldoubleturn \quad \mathcal{Q}
	\qquad \text{ ssi } \qquad 
	\mathcal{T} \quad \lturn \quad \mathcal{Q}
	$}	
\end{theoreme}

Reformulation :
\begin{theoreme}[Théorème B] 
	\sauteligne
	\mybox{$\mathcal{T}$ cohérente  
		\qquad ssi \qquad 
		$\mathcal{T}$ admet un modèle
	}	
\end{theoreme}

On pourrait énoncer un résultat similaire pour un ensemble quelconque de formules :
\mycenterline{
$\Gamma$ cohérent  
\qquad ssi \qquad 
$\Gamma$ satisfaisable.	
}

%--------------------------------------------------------------------
\subsection{Démonstration du théorème B à partir du théorème A}

\begin{align*}
	&  \mathcal{T} \ \text{ non cohérente}\\
	\text{ ssi } \quad 	& \mathcal{T} \ \ \lturn \ \ \lantitauto \qquad \text{par définition}\\
	\text{ ssi } \quad 	&  \mathcal{T} \ \ \ldoubleturn \ \ \lantitauto  \qquad \text{par le théorème A}\\
	\text{ ssi } \quad 	&  \mathcal{T} \ \text{n'est satisfaite par aucune structure $\mathcal{S}$}\\	
	\text{ ssi } \quad 	&  \mathcal{T} \ \text{n'a pas de modèle}
\end{align*} 

On verra en exercice (et dans le prochain chapitre) qu'on pourrait prouver le théorème A à partir du théorème B.



%%%%%%%%%%%%%%%%%%%%%%%%%%%%%%%%%%%%%%%%%%%%%%%%%%%%%%%%%%%%%%%%%%%%%
\section{Application à la finitude}

Nous nous intéressons au cas où l'ensemble des hypothèses est infini (mais dénombrable).
Nous allons voir que seul un nombre fini d'hypothèses est nécessaire.
Nous aurons deux résultats, un sur les preuves (le théorème de finitude) et un sur la satisfaction (le théorème de compacité en logique du premier ordre, analogue à celui du chapitre \og{}Théorème de compacité\fg{} de la logique V/F).


%--------------------------------------------------------------------
\subsection{Théorème de finitude}

Soit $\Gamma$ un ensemble de formules (éventuellement infini), soit $\mathcal{Q}$ une formule d'un langage $\mathcal{L}$.

\begin{theoreme}[Théorème de finitude]
	Si $\Gamma \ \lturn \ \mathcal{Q}$ alors il existe une partie finie $\Gamma_0 \subset \Gamma$ telle que $\Gamma_0 \ \lturn \ \mathcal{Q}$.
\end{theoreme}
\index{theoreme@théorème!de finitude}


Remarque : la réciproque est bien sûr vraie.

\begin{proof}
La démonstration est fondée sur la nature même de ce qu'est une preuve.
$\Gamma \ \lturn \ \mathcal{Q}$ signifie qu'il existe une preuve de $\mathcal{Q}$ à partir des formules de $\Gamma$. Considérons une telle preuve.
Cette preuve consiste en un nombre fini d'étapes et à chaque étape on applique une règle. Chaque règle ne fait intervenir qu'un nombre fini de formules de $\Gamma$.
On note $\Gamma_0$ l'ensemble des formules de $\Gamma$ qui apparaissent dans cette preuve. Alors on a bien obtenu une preuve de $\mathcal{Q}$ à partir de $\Gamma_0$, donc $\Gamma_0 \ \lturn \ \mathcal{Q}$ avec $\Gamma_0$ fini.	
\end{proof}

%--------------------------------------------------------------------
\subsection{Théorème de compacité}

La version correspondante pour la satisfaction est le théorème de compacité pour la logique du premier ordre, similaire à celui rencontré pour la logique V/F dans le chapitre \og{}Théorème de compacité\fg{} de la première partie.


\begin{theoreme}[Théorème de compacité]
	$\Gamma$ satisfaisable ssi toute partie finie $\Gamma_0 \subset \Gamma$ est satisfaisable.
\end{theoreme}
\index{theoreme@théorème!de compacite@de compacité}

Reformulation : $\Gamma \ \ldoubleturn \ \mathcal{Q}$ ssi pour toute partie finie $\Gamma_0 \subset \Gamma$ on a $\Gamma_0 \ \ldoubleturn \ \mathcal{Q}$. 

Remarque : à chaque fois le sens direct $\Rightarrow$ est évident.

Pour les théories, une reformulation du sens intéressant est :
\begin{theoreme}
	Si toute partie finie d'une théorie $\mathcal{T}$ admet un modèle, alors $\mathcal{T}$ elle-même admet un modèle.
\end{theoreme}

Il est intéressant de prendre le temps de réfléchir à la profondeur d'un tel théorème.
Chaque partie finie a un modèle, mais bien sûr ce modèle peut être différent pour chaque $\Gamma_0$.
Cependant, on trouve à la fin un modèle commun à toutes les formules.

\begin{proof}[Démonstration du théorème de compacité]
On a déjà dit que le sens direct $\Rightarrow$ est évident.
On va démontrer le sens réciproque $\Leftarrow$ par contraposition : \og{}si $\Gamma$ non satisfaisable, alors il existe $\Gamma_0$ fini tel que $\Gamma_0$ non satisfaisable.\fg{}	
	
\begin{align*}
	&  \Gamma \ \text{ non satisfaisable}\\
	\text{ donc } \quad 	& \Gamma \ \ \ldoubleturn \ \ \lantitauto \qquad \text{car de toute façon $\Gamma$ n'est jamais satisfait}\\
	\text{ donc } \quad 	& \Gamma \ \ \lturn \ \ \lantitauto \qquad \text{par le théorème de complétude, donc $\Gamma$ non cohérente}\\
	\text{ donc } \quad 	& \Gamma_0 \ \ \lturn \ \ \lantitauto \qquad \text{pour un certain $\Gamma_0 \subset \Gamma$ fini, par le théorème de finitude}\\
	\text{ donc } \quad 	& \Gamma_0 \ \ \ldoubleturn \ \ \lantitauto \qquad \text{par le théorème de validité}\\			
	\text{ donc } \quad 	& \Gamma_0 \ \text{non satisfaisable}\\
	\text{ donc } \quad 	& \text{il existe $\Gamma_0$ fini, non satisfaisable}\\			
	\text{ donc } \quad 	& \text{non\;(tout $\Gamma_0$ fini est satisfaisable)}		
\end{align*} 	
	
\end{proof}

%%%%%%%%%%%%%%%%%%%%%%%%%%%%%%%%%%%%%%%%%%%%%%%%%%%%%%%%%%%%%%%%%%%%%
\section{Théorème de clôture}


Le théorème de clôture explique que l'on peut passer d'un énoncé du théorème de complétude avec des formules fermées à un énoncé pour des formules quelconques.

%--------------------------------------------------------------------
\subsection{Énoncé}



\begin{theoreme}[Théorème A] 
Soit $\mathcal{T}$ une théorie, $\mathcal{Q}$ une formule fermée.
	\[	
		\text{Si } \ \mathcal{T} \quad \ldoubleturn \quad \mathcal{Q}
		\qquad \text{ alors } \qquad 
		\mathcal{T} \quad \lturn \quad \mathcal{Q}.
   \]	
\end{theoreme}


\begin{theoreme}[Théorème C] 
Soit maintenant $\Gamma$ un ensemble de formules quelconques, et $\mathcal{Q}$ une formule quelconque.
	\[	
	\text{Si } \ \Gamma \quad \ldoubleturn \quad \mathcal{Q}
	\qquad \text{ alors } \qquad 
	\Gamma \quad \lturn \quad \mathcal{Q}.
	\]	
\end{theoreme}

\begin{theoreme}[Théorème de clôture]
	Le théorème A entraîne le théorème C.
\end{theoreme}
\index{theoreme@théorème!de cloture@de clôture}

Autrement dit, pour démontrer le théorème de complétude en toute généralité, il suffit de démontrer le théorème de complétude pour les formules fermées.

%--------------------------------------------------------------------
\subsection{Clôture}

Avant de donner la démonstration du théorème C, introduisons une notion de clôture par introduction de nouvelles constantes.

\begin{definition}
	Soit $\mathcal{P}$ une formule ayant des variables libres, notées $x_1, \ldots, x_n$, définie sur un langage $\mathcal{L}$.
	On note $\mathcal{L}^c$ le langage obtenu à partir de $\mathcal{L}$ en ajoutant $n$ nouvelles constantes $c_1, \ldots, c_n$ (ces constantes ne doivent pas apparaître dans $\mathcal{L}$).
	On note $\mathcal{P}^c$ \defi{substitution par les constantes nouvelles} de $\mathcal{P}$, la formule obtenue à partir de $\mathcal{P}$ en remplaçant chaque variable libre $x_i$ par la nouvelle constante $c_i$.
\end{definition}
Par construction, $\mathcal{P}^c$ est une formule fermée (du langage $\mathcal{L}^c$).

Pour un ensemble de formules $\Gamma$ et une formule $\mathcal{Q}$, on note $x_1, \ldots, x_n$ l'ensemble des variables libres dans $\Gamma \cup \{ \mathcal{Q} \}$.
On construit ainsi un langage $\mathcal{L}^c$, et des formules fermées $\Gamma^c$ (la clôture par les constantes nouvelles des formules de $\Gamma$) et $\mathcal{Q}^c$.

Voici quelques propriétés utiles pour la preuve, qui sont essentiellement des jeux d'écriture.
\begin{itemize}
	\item Point a.
	
	\[
	\Gamma \ \ldoubleturn  \ \mathcal{Q} \qquad \text{ implique } \qquad \Gamma^c \ \ldoubleturn \ \mathcal{Q}^c
	\]


	Justifions en détail. 
	Considérons une structure $\mathcal{S}^c$ (pour le langage $\mathcal{L}^c$) qui satisfait $\Gamma^c$. 
	On note $d_i = c_i^{\mathcal{S}^c}$ l'interprétation de la constante $c_i$ dans $\mathcal{S}^c$.
	
	On peut construire une structure $\mathcal{S}$ (pour le langage $\mathcal{L}$) en oubliant les constantes nouvelles $c_1, \ldots, c_n$ et en même temps on construit une affectation $s$ en posant $s(x_i) = d_i$ pour chaque variable libre $x_i$. Pour une formule $\mathcal{P}$ (de $\Gamma$ ou $\mathcal{Q}$), on a alors :
	\[
	\ldoubleturn_{\mathcal{S}^c} \mathcal{P}^c  \qquad \text{ ssi } \qquad \ldoubleturn_{\mathcal{S}, s} \mathcal{P}
	\]
	(Ceci découle directement de la substitution de chaque variable libre $x_i$ par la constante $c_i$.)
	
	Comme $\ldoubleturn_{\mathcal{S}^c} \Gamma^c$, alors on a $\ldoubleturn_{\mathcal{S}, s} \Gamma$ (par la substitution précédente). 
	Par hypothèse, on a $\Gamma \ \ldoubleturn  \ \mathcal{Q}$, donc
	$\ldoubleturn_{\mathcal{S}, s} \mathcal{Q}$, et donc en retour $\ldoubleturn_{\mathcal{S}^c} \mathcal{Q}^c$.

	\item Point b. Le lemme des constantes nouvelles :
	\[
	\Gamma^c \ \lturn  \ \mathcal{Q}^c \qquad \text{ implique } \qquad \Gamma \ \lturn \ \mathcal{Q}
	\]	
	
	En effet imaginons qu'on a une preuve écrite en déduction naturelle de $\mathcal{Q}^c$ à partir de $\Gamma^c$.
	On recopie cette preuve, en remplaçant dans toutes les lignes les occurrences des constantes nouvelles $c_i$ par les variables $x_i$. On obtient une preuve de $\mathcal{Q}$ à partir de $\Gamma$.

\noindent
\begin{minipage}[t]{0.4\textwidth}
    \centering
    \textbf{Preuve de $\Gamma^c \ \lturn \ \mathcal{Q}^c$}
    \[
    \begin{nd}
        \hypo [~] {} {\mathcal{P}_1^c} \by{\text{Hypothèses $\Gamma^c$}}{}
      	\hypo [~] {} {\vdots}		
        \hypo [~] {} {\mathcal{P}_l^c}
        \have [~] {} {\vdots} \by{\text{Preuve}}{}
		\have [~] {} {\ldots c_i \ldots} \by{\text{Les constantes nouvelles}}{}
        \have [~] {} {\vdots}		
        \have [\lturn] {} {\mathcal{Q}^c} \by{\text{Conclusion}}{}
    \end{nd}
    \]
\end{minipage}
\qquad\qquad
\begin{minipage}[t]{0.4\textwidth}
    \centering
    \textbf{Preuve de $\Gamma \ \lturn \ \mathcal{Q}$}
    \[
    \begin{nd}
        \hypo [~] {} {\mathcal{P}_1} \by{\text{Hypothèses $\Gamma$}}{}
      	\hypo [~] {} {\vdots}
        \hypo [~] {} {\mathcal{P}_l}
        \have [~] {} {\vdots} \by{\text{Preuve}}{}
		\have [~] {} {\ldots x_i \ldots} \by{\text{Retour aux variables}}{}
        \have [~] {} {\vdots}	
        \have [\lturn] {} {\mathcal{Q}} \by{\text{Conclusion}}{}
    \end{nd}
    \]
\end{minipage}	
	
\end{itemize}

%--------------------------------------------------------------------
\subsection{Démonstration du théorème C à partir du théorème A}

\begin{align*}
						& \Gamma \ \ldoubleturn  \ \mathcal{Q} \\
	\text{ donc } \quad & \Gamma^c	\ \ldoubleturn \ \mathcal{Q}^c \qquad \text{par le point a}\\
	\text{ donc } \quad & \Gamma^c	\ \lturn \ \mathcal{Q}^c \qquad \text{par le théorème A pour les formules fermées}\\
	\text{ donc } \quad & \Gamma \ \lturn \ \mathcal{Q} \qquad \text{par le point b}	
\end{align*}


%--------------------------------------------------------------------
\subsection{Clôture universelle}

\index{cloture universelle@clôture universelle}

Il existe une autre notion de clôture, appelée \emph{clôture universelle}.

\begin{definition}
	Soit $\mathcal{P}$ une formule ayant des variables libres, notées $x_1, \ldots, x_n$. On note $\overline{\mathcal{P}}$ (ou aussi $\forall \mathcal{P}$) la \defi{clôture universelle} de $\mathcal{P}$ :
	\[
	\overline{\mathcal{P}} : \qquad 
	\forall x_1 \ \forall x_2 \ \ldots \ \forall x_n \quad \mathcal{P}
	\] 
\end{definition}




%%%%%%%%%%%%%%%%%%%%%%%%%%%%%%%%%%%%%%%%%%%%%%%%%%%%%%%%%%%%%%%%%%%%%
\section*{Exercices}

\subsection*{Théorème de complétude}

\begin{exercicecours}
On considère $\Gamma$ un ensemble de formules et $\mathcal{P}$, $\mathcal{Q}$ des formules telles que :
\begin{equation}
	\label{exo:completude:1}
	\tag{$\star$}
\Gamma \ \  \lturn \ \  \mathcal{P} 
\qquad \text{ et } \qquad
\mathcal{P} \ \ \lturn \ \ \mathcal{Q} 
\end{equation}
Le but de l'exercice est d'obtenir $\Gamma \ \lturn \ \mathcal{Q}$.

\begin{enumerate}
	\item \emph{Première méthode.}
	À partir d'une preuve en déduction naturelle de $\Gamma \ \lturn \ \mathcal{P}$ et d'une preuve  en déduction naturelle de $\mathcal{P} \ \lturn \ \mathcal{Q}$, écrire une preuve de $\Gamma \ \lturn \ \mathcal{Q}$. 
	
	\item\emph{Deuxième méthode.}
	\begin{enumerate}
		\item Quels énoncés obtient-on à partir de nos hypothèses \eqref{exo:completude:1} en appliquant le théorème de validité ?
		
		\item Sous les hypothèses \eqref{exo:completude:1}, montrer qu'on a $\Gamma \ \ldoubleturn \mathcal{Q}$.
		
		\item Conclure.
	\end{enumerate}	
\end{enumerate}	
\end{exercicecours}


\begin{exercicecours}
Soit $\mathcal{L}$ un langage muni d'un prédicat binaire $R(x,y)$.
On considère les trois formule suivantes :
\begin{enumerate}
	\item symétrie $\mathcal{S} : \forall x \ \forall y \ \ \big( R(x,y) \limplies R(y,x) \big)$
	
	\item transitivité $\mathcal{T} : \forall x \ \forall y \ \forall z \ \ \big( R(x,y) \land R(y,z) \limplies R(x,z) \big)$	
	
	\item réflexivité $\mathcal{Q} : \forall x \  \ R(x,x)$
\end{enumerate}
Une relation $R$ pour laquelle ces trois formules sont vraies, s'appelle une \emph{relation d'équivalence}. On se demande si les deux seuls axiomes $\mathcal{S}$ et $\mathcal{T}$ suffisent, c'est-à-dire a-t-on $\mathcal{S}, \ \mathcal{T} \ \lturn \ \mathcal{Q}$ ?

\begin{enumerate}
	\item Un étudiant propose la preuve suivante :
	\begin{quote}
	\emph{Soit $x$ un élément quelconque. Prenons un élément $y$ tel que $R(x,y)$. Par l'axiome $\mathcal{S}$ (symétrie), on a aussi $R(y,x)$. En appliquant l'axiome $\mathcal{T}$ (transitivité) sur $R(x,y)$ et $R(y,x)$, on en déduit $R(x,x)$. La relation est donc réflexive	et $\mathcal{Q}$ est prouvée.}
	\end{quote}
	
	Quelle est l'erreur logique fondamentale dans ce raisonnement ?
	
	\item Trouver une structure $\mathcal{S}$, c'est-à-dire définir un domaine $E$ et y expliciter la relation $R$, telle que les formules  $\mathcal{S}$ et $\mathcal{T}$ y soient vraies, mais telle que $\mathcal{Q}$ ne le soit pas.
	
	\item A-t-on $\mathcal{S}, \ \mathcal{T} \ \lturn \ \mathcal{Q}$ ?
\end{enumerate}
\end{exercicecours}	
	
\subsection*{Applications aux théories}

\begin{exercicecours}
	Montrer l'équivalence entre ces deux définitions d'une théorie $\mathcal{T}$ \emph{non cohérente} :
	\begin{itemize}
		\item[(i)] il existe $\mathcal{Q}$, telle que  $\mathcal{T} \ \lturn \ \mathcal{Q}$ et $\mathcal{T} \ \lturn \ \lnot \mathcal{Q}$,
		
		\item[(ii)] $\mathcal{T} \ \lturn \ \lantitauto$.
	\end{itemize}
\end{exercicecours}


\begin{exercicecours}
	À l'aide du théorème de complétude \og{}si $\mathcal{T} \ \ldoubleturn \ \mathcal{Q}$ alors 
    $\mathcal{T} \ \lturn \ \mathcal{Q}$\fg{} prouver que :
    \mycenterline{
    \og{}si $\mathcal{T}$ est insatisfaisable alors $\mathcal{T}$ est non cohérente\fg{}. 	
    }	
    Démontrer ensuite la réciproque.
\end{exercicecours}


\begin{exercicecours}
	On rappelle les théorèmes A et B :
	\begin{itemize}
		\item Théorème A. 
		\[	
		\mathcal{T} \quad \ldoubleturn \quad \mathcal{Q}
		\qquad \text{ ssi } \qquad 
		\mathcal{T} \quad \lturn \quad \mathcal{Q}.
		\]
		
		\item Théorème B.
		\[
		\mathcal{T} \text{ cohérente }
		\qquad \text{ ssi }\qquad 
		\mathcal{T} \text{ admet un modèle}.
		\]
	\end{itemize}
	
	On admet le théorème B. En déduire le théorème A.	
\end{exercicecours}



\subsection*{Application à la finitude}

\begin{exercicecours}
	Soit $\mathcal{L}$ un langage (contenant le symbole d'égalité $=$). 
	Pour tout entier $n \geq 1$, on note $\mathcal{P}_n$ la formule exprimant \og{}il existe au moins $n$ éléments distincts \fg{}.
	
	\begin{enumerate}
		\item Écrire explicitement les formules $\mathcal{P}_1$, $\mathcal{P}_2$, $\mathcal{P}_3$.
		
		\item Soit $\mathcal{T}$ une théorie ayant des modèles finis de taille arbitrairement grande, c'est-à-dire que pour tout entier $k$, $\mathcal{T}$ admet un modèle ayant au moins $k$ éléments. 
		
		On pose $\Gamma = \mathcal{T} \cup \{ \mathcal{P}_n \mid n \geq 1 \}$. Montrer que tout sous-ensemble fini de $\Gamma$ admet un modèle.
		
		\item En utilisant le théorème de compacité, en déduire que $\mathcal{T}$ admet un modèle infini.
		
		\item Conclusion : existe-t-il une théorie du premier ordre $\mathcal{T}_{\text{fini}}$ dont les modèles seraient \textbf{exactement} toutes les structures finies ?
	\end{enumerate}
\end{exercicecours}


\subsection*{Théorème de clôture}

% \begin{exercicecours}
% Voici une démonstration alternative du théorème de clôture.	
% Compléter les termes manquants et donner une justification pour le passage de chaque ligne à la suivante.
% \begin{align*}
% 						& \Gamma \ \ldoubleturn  \ \mathcal{Q} \\
% 	\text{ donc } \quad & \overline{\Gamma}	\ \ldoubleturn \ \overline{\mathcal{Q}} \qquad \text{d'après \ldots}\\
% 	\text{ donc } \quad & \overline{\Gamma}	\cup \{ \lnot \overline{\mathcal{Q}} \}  \text {\  \ est \ldots } \\
% 	\text{ donc } \quad & \overline{\Gamma}	\cup \{ \lnot \overline{\mathcal{Q}} \}  \text {\  \ est \ldots } \\
% 	\text{ donc } \quad & \overline{\Gamma}	\cup \{ \lnot \overline{\mathcal{Q}} \} \ \lturn \ \lantitauto \qquad \text { \ldots } \\		
% 	\text{ donc } \quad & \overline{\Gamma}	\ \lturn \ \overline{\mathcal{Q}} \qquad \text{car \ldots }\\	
% 	\text{ donc } \quad & \overline{\Gamma}	\ \lturn \ \mathcal{Q} \qquad \text{car \ldots }\\
% 	\text{ donc } \quad & \Gamma \ \lturn \ \mathcal{Q} \qquad \text{ car \ldots }
% \end{align*}
% \end{exercicecours}


\begin{exercicecours}
	Pour une formule $\mathcal{P}$ dont on note $x_1$, \ldots, $x_n$ les variables libres, sa \emph{clôture existentielle} est la formule $\mathcal{P}_{\exists}$ (notée aussi $\exists \mathcal{P}$) :
		\[
	\mathcal{P}_{\exists} : \qquad 
	\exists x_1 \ \ldots \ \exists x_n \quad \mathcal{P}
	\] 
	\begin{enumerate}
		\item Montrer que $\mathcal{P}$ est satisfiable si et seulement si $\mathcal{P}_{\exists}$ admet un modèle.
		
		
		\item Montrer l'équivalence $\lnot \big( \overline{\mathcal{P}} \big) \ \equiv \ (\lnot \mathcal{P})_{\exists}$. 
		
		On rappelle que $\overline{\mathcal{P}}$ désigne la clôture universelle et que $\mathcal{P} \equiv \mathcal{Q}$ signifie ($\mathcal{P} \ldoubleturn \mathcal{Q}$ et $\mathcal{Q} \ldoubleturn \mathcal{P}$), voir le chapitre \og{}Satisfaction\fg{}.
		
		\item Montrer que si $\Gamma \ \lturn \ \mathcal{Q}$, alors $\Gamma \ \lturn \mathcal{Q}_{\exists}$.
		La réciproque est-elle vraie ?
	\end{enumerate}
\end{exercicecours}	
	
\end{document}
